\scanpage{84}
\margin{$\unsure{\epsilon|E|c}$}

$\mathcal M$ catégorie motivique \inserted{$1/2$ simple}, de gerbe $\mathscr G$. Supposons que le lien $\mathcal L$ de $\mathscr G$ soit \uline{commutatif}, donc correspond à un groupe algébrique \inserted{comm de type multiplicatif} $G$ sur $\mathbb Q$. Ainsi, $G^0$ est un tore. Soit $E$ un motif ``fidèle''. \deleted{et soit $F$ un foncteur} \inserted{Alors $G$ opère} sur $E$, considéré comme \deleted{objet} \inserted{section} de la catégorie fibrée $\mathrm{HOM}(\mathscr G,\underline{\underline{\unsure{\mathrm{Proj}|\mathrm{Rep}|\mathrm{Prj}}}})$, et en particulier $G(\mathbb Q)$ opère sur $E$ donc $G(\mathbb Q)\overset{\varphi}{\hookrightarrow}\mathrm{End}_{\mathcal M}(E)$. D'ailleurs,

\deleted{les \unsure{élémts} de $G(\mathbb Q)$}

\deleted{Soit $U=\mathrm{Ker}(G^0\longrightarrow G_m)$, alors $U$ opère}

\deleted{\unsure{La}} Choisissons une forme de polarisation sur $E$, d'où une structure d'algèbre de Riemann sur $A=\mathrm{End}_{\mathcal M}(E)$, et on trouve que pour $u\in G(\mathbb Q)$, $\varphi(u)\varphi(u)^{\deleted{\unsure{*}}'}=\varepsilon(u)1\deleted{\unsure{*|E}}$.

[En particulier, si $u\in\inserted{U=}\mathrm{Ker}\,\varepsilon$, on aura $\varphi(u)\varphi(u)'=\mathrm{id}$.]. Donc l'algèbre engendrée par les $\varphi(u)$ dans $A$ est auto-adjointe, et $U(\mathbb Q)$ apparaît comme un ss-groupe du ``groupe unitaire'' de $A$, formé des $u$ tels que $uu'=1$. \deleted{Plus précisément} Or ce \inserted{dernier} groupe est \inserted{(formé sur $\unsure{\mathrm{Vect}.|\mathrm{Vert}.|\mathrm{vect}}/\mathbb Q$) et} un \uline{tore} \uline{compact}. On trouve que $U^0$ est un \uline{tore compact}.
% archive sheet number: 83

